\documentclass[11pt,letterpaper]{article}
\usepackage[margin=0.85in,headheight=16pt]{geometry}
\usepackage[T1]{fontenc}
\usepackage[utf8]{inputenc}
\usepackage{lmodern,microtype}
\usepackage{amsmath,amssymb,mathtools,booktabs,longtable,tabularx,array}
\usepackage[dvipsnames]{xcolor}
\usepackage{hyperref,fancyhdr,titlesec,enumitem,graphicx,tcolorbox}
\usepackage{lastpage}
\definecolor{ink}{HTML}{14263B}
\definecolor{sea}{HTML}{167C8A}
\definecolor{gold}{HTML}{C49747}
\definecolor{pale}{HTML}{EEF5F6}
\definecolor{graytext}{HTML}{516275}
\hypersetup{colorlinks=true,linkcolor=ink,urlcolor=sea,bookmarksopen=true,pdfauthor={Independent Educational Engineering Suite}}
\titleformat{\section}{\Large\bfseries\color{ink}}{\thesection}{0.7em}{}[\vspace{-0.3em}\color{gold}\titlerule]
\titleformat{\subsection}{\large\bfseries\color{sea}}{\thesubsection}{0.6em}{}
\setlist[itemize]{leftmargin=1.3em,itemsep=2pt,topsep=3pt}
\setlist[enumerate]{leftmargin=1.5em,itemsep=2pt,topsep=3pt}
\pagestyle{fancy}
\fancyhf{}
\fancyhead[L]{\small\sffamily\color{graytext}SECURITY ANALYSIS / ENGINEERING SUITE}
\fancyhead[R]{\small\sffamily\color{graytext}\textbf{\thepage}/\pageref{LastPage}}
\renewcommand{\headrulewidth}{0.3pt}
\renewcommand{\footrulewidth}{0pt}
\newtcolorbox{methodbox}[1][]{colback=pale,colframe=sea,boxrule=0.55pt,arc=1mm,title={\bfseries Method note},fonttitle=\sffamily,#1}
\newtcolorbox{sourcebox}[1][]{colback=white,colframe=gold,boxrule=0.55pt,arc=1mm,title={\bfseries Source boundary},fonttitle=\sffamily,#1}
\newcommand{\term}[1]{\textbf{\color{ink}#1}}
\newcommand{\src}{\textsf{\color{sea}[SOURCE]}}
\newcommand{\ext}{\textsf{\color{gold}[EXTENSION]}}
\newcommand{\gap}{\textsf{\color{red}[SOURCE GAP]}}
\newcommand{\answer}[1]{\par\smallskip\noindent\textbf{Check:} #1\par}
\newcommand{\cover}[3]{%
\begin{titlepage}\pagecolor{ink}\color{white}\vspace*{1.15in}
{\sffamily\large\color{gold}INDEPENDENT TECHNICAL DOCUMENTATION SUITE\par}\vspace{0.4in}
{\Huge\sffamily\bfseries #1\par}\vspace{0.15in}
{\Large\sffamily #2\par}\vspace{0.65in}
\color{gold}\rule{\textwidth}{2pt}\par\vspace{0.4in}
{\large #3\par}\vfill
{\small\sffamily Based on a critical reading of Graham and Dodd's \emph{Security Analysis}, Seventh Edition (2023). Original educational analysis; not affiliated with the publisher.\par}
\end{titlepage}\nopagecolor\color{black}}
\newcommand{\front}[3]{\cover{#1}{#2}{#3}\tableofcontents\clearpage
\begin{sourcebox}\src\ means a concept is attributable to the relevant chapter or seventh-edition contributor, summarized independently rather than reproduced. \ext\ means a modern method or engineering recommendation prepared for this suite, not an original Graham--Dodd rule. Historical examples and standards should not be treated as current market requirements. All numerical examples are synthetic, educational, and not securities recommendations.\end{sourcebox}}
\setlength{\parskip}{5pt}\setlength{\parindent}{0pt}
\begin{document}
\front{Security Analysis}{Technical Reference and Analytical Architecture}{Part-aligned terminology, valuation foundations, claim analysis, accounting and controls}

\section{Analytical charter and investment decision}
\src\ Chapters 1--5 distinguish an evidence-based security analysis from mere quotations or speculation. Intrinsic value is a reasoned appraisal subject to error; the analyst must also study the instrument's contractual claim. A stock is a residual economic ownership interest, while a bond is ordinarily a contractual senior claim. Neither label guarantees safety.

\ext\ \term{Working decision object.} Record a security identifier, effective date, legal issuer, currency, adjusted financials, claim ranking, valuation range, downside state, catalyst, liquidity, suitability, and approval trail. Use a thesis whose falsification criteria are written before position initiation.

\begin{methodbox}
\textbf{Four independent gates:} (1) evidentiary sufficiency, (2) contractual and financial survival, (3) valuation with margin of safety, and (4) portfolio fit. An attractive expected return does not override a failed solvency or mandate gate.
\end{methodbox}

\subsection{Intrinsic value and discount discipline}
Define price $P$, analyst value interval $[V_L,V_H]$, and a selected underwriting point estimate $V_B$. When $V_B>0$,
\[
\operatorname{MOS}_{B}=1-\frac{P}{V_B},\qquad
\operatorname{MOS}_{L}=1-\frac{P}{V_L}.
\]
The second expression is a stronger stress test, not necessarily a statistical confidence bound. Market price can rationally diverge from business value due to forced trades, information disagreement, constraints, illiquidity, or a genuinely deteriorating business. A positive calculated discount is not sufficient evidence of a bargain.

\subsection{Evidence and claim hierarchy}
\begin{center}\small
\begin{tabularx}{\textwidth}{p{0.2\textwidth}X X}
\toprule Evidence type & Examples & Verification/control \\
\midrule Primary filing & Audited statements, notes, debt schedules & Entity/date, auditor opinion, restatement tracking \\
Contractual & Indentures, guarantee and security agreements & Claim rank, baskets, triggers, release provisions \\
Operational & Unit economics, customers, operating metrics & Cross-check against cash conversion \\
Independent & Competitor filings, rate cases, industry datasets & Disagreement log and range of outcomes \\
Management claim & Guidance and adjusted EBITDA & Reconcile to GAAP/IFRS and challenge incentives \\
\bottomrule
\end{tabularx}
\end{center}

\section{Fixed-value investments and creditor analysis}
\src\ Parts II and III (Chapters 6--26) emphasize adverse-condition solvency, relative security of contractual income and principal, enforceable protective covenants, seniority, collateral, guarantees, preferred terms, conversion rights, and questionable-safety claims.

\subsection{Coverage, leverage and payment capacity}
For interest-bearing obligations, define interest coverage $IC=\text{EBIT}/I$ for interest expense $I>0$. A useful credit interrogation pairs historical trough coverage with forward stress metrics and cash-flow liquidity. It is dangerous to use a single coverage ratio as the sole decision rule. When $I\leq 0$, label the metric not applicable instead of manufacturing an infinite ratio.

\ext\ Build a debt maturity ladder, available revolver capacity, unrestricted cash by legal entity, covenant-headroom worksheet and 12--24 month stressed liquidity runway. Use both operating cash flow after necessary reinvestment and debt-service obligations. Interest coverage, fixed-charge coverage, and debt service coverage differ in numerator, denominator, and appropriateness to the issuer.

\subsection{Claims waterfall and structural subordination}
At enterprise realizable proceeds $A$ net of court/transaction costs, distribute value according to enforceable priority, available collateral, intercreditor contracts, and jurisdictional law. Illustrative simplified debt recovery, \emph{not a legal prediction}, is
\[
R_j=\frac{\min\bigl(C_j,\max(0,A-\sum_{k<j} C_k)\bigr)}{C_j},
\]
where $C_j$ is claim $j$ and earlier $k$ represent fully prior claims in the same simplified waterfall. Replace this toy model with a legal-entity and collateral-specific model in actual underwriting. Account for administrative expenses, taxes, DIP financing, cross guarantees, and avoidance risks.

\subsection{Optionality in senior instruments}
\src\ Chapters 22--25 analyze participating senior securities, conversion mechanics, warrants, and switching/hedging. \ext\ For a simplified convertible, compute conversion shares $q=F/K$ for face value $F$ and conversion price $K$. Conversion parity is $qS$ and premium is $P_c/(qS)-1$ when $qS>0$. The convertible has contingent value from both debt claim and equity option; the intuitive ``bond floor + option'' decomposition requires care because the two are not independent under credit deterioration.

\section{Common-stock earning power and distributions}
\src\ Chapters 27--30 distinguish speculative growth hopes from repeatable earning power, dividend records and distributable business value. Modern analysis should preserve the skepticism of unexamined assumptions while recognizing that sustainable growth can affect fundamental worth.

\subsection{Normalization and owner earnings}
\ext\ Analytical normalization starts with reported operating results, reverses only demonstrably unusual items, and estimates sustainable reinvestment needs. Free cash flow to the firm (FCFF) is
\[
FCFF=EBIT(1-T)+D\!\&A-CapEx-\Delta NWC,
\]
with tax rate $T$, net capital spending and operating net-working-capital investment defined consistently. FCFF is not equal to EBITDA and not equal to unrestricted cash available to shareholders. Disaggregate acquired earnings, organic growth, price/volume/mix, and changes in economic working capital.

\subsection{Dividend and repurchase discipline}
A dividend is economically funded only when cash generation, balance-sheet flexibility, and reinvestment needs permit it. Track recurring payout ratios and debt-funded distributions; for repurchases, track price versus value, dilution from compensation, changes in shares outstanding, and alternative uses of capital.

\section{Income-account integrity and capitalization}
\src\ Part V (Chapters 31--41) examines extraordinary losses, potentially misleading reported income, depreciation, earning records, price--earnings ratios, capitalization, and source of income. Earnings need to be recast for repeatability and economic interpretation.

\subsection{Normalization controls}
\begin{enumerate}
\item Preserve the original filing figures before adjusting.
\item Identify whether an item is unusual in nature \emph{and} unlikely to recur; classify repeat restructurings as recurring where indicated by evidence.
\item Reconcile noncash charges to maintenance capital and economic impairment.
\item Reconcile consolidated earnings to parent-available distributable cash.
\item State after-tax impact, per-share impact, and source for every adjustment.
\end{enumerate}

\subsection{Capitalization bridge}
\ext\ Under a consistent enterprise-value definition,
\[
\text{Equity value}=EV+\text{nonoperating cash}-\text{debt}-\text{preferred claims}-\text{other senior claims}.
\]
Noncontrolling interests and associates require explicit treatment, depending on how operating metrics are consolidated. Fully diluted per-share value uses treasury-stock or if-converted mechanics appropriate to the instrument, rather than simply counting all options at face value.

\section{Balance-sheet value and liquidation}
\src\ Chapters 42--45 evaluate book value, current asset value, stockholder--management relations, liquidation prospects, and hidden balance-sheet weakness. The historic net-current-asset concept is a protective asset measure, not a universal proxy for going-concern value.

\subsection{NCAV and conservative asset realization}
\[
NCAV=\text{current assets}-\text{total liabilities},\qquad
\text{NCAV/share}=NCAV/N.
\]
Positive NCAV does not mean current assets can be converted to cash without losses. \ext\ Haircut receivables and inventory; include obligations omitted from simple statement categories. For liquidation, estimate $L=(\sum_i \alpha_i A_i)-\sum_j C_j-\text{sale expenses}$ with documented asset recovery fractions $0\leq\alpha_i\leq1$ and all relevant claims.

\section{Relative value, governance and market discrepancies}
\src\ Chapters 46--52 cover warrants, financing cost, pyramids, peer comparisons, gaps between price and value, and the distinction between market analysis and security analysis. A low multiple without peer and governance adjustment is an unreliable value signal.

\ext\ Maintain a peer-normalization ledger (fiscal year, lease policy, capital intensity, adjusted operating margin, ROIC, currency, leverage) and compare valuations both before and after quality adjustment. For corporate control pyramids, separately diagram voting rights, economic participation, and debt located at each layer. Include governance, holding-company discounts and uncertainty in realization timing.

\section{Seventh-edition perspectives and international investing}
\src\ Part VIII consists of contributions on international investing, public/private investing, and endowment management. These chapters add perspective on currency and accounting reliability, investment time horizon, control and liquidity differences, and institutional constraints. The seventh-edition preface also emphasizes uncertainty, downside protection, judgment, and future cash flows alongside the classic framework.

\ext\ Record home and local currencies, translation and convertibility restrictions, sovereign and transfer risks, controlling-shareholder rights, withholding tax, trading liquidity, and custody. Separate private-market appraisal uncertainty from the genuine option to wait, which may have value only for unforced capital. Endowment investment policy should specify liquidity budget and spending obligations before pursuing illiquid return premiums.

\section{Reference glossary and crosswalk}
\small
\begin{longtable}{p{0.24\textwidth}p{0.33\textwidth}p{0.34\textwidth}}
\toprule Term & Analytic meaning & Misuse to avoid \\
\midrule\endhead
Intrinsic value & Business/claim worth inferred from fundamentals and enforceable rights & Treating a quote or single DCF output as truth \\
Margin of safety & Discount intended to accommodate error and adversity & Guarantee of no loss \\
Earning power & Maintainable earning capacity over representative conditions & Latest-quarter extrapolation \\
Interest coverage & Selected earnings/cash numerator relative to interest obligations & Replacing liquidity and maturity analysis \\
Protective covenant & Contractual restriction or remedy for creditors & Assuming breach leads to automatic full recovery \\
Preferred stock & Contract-defined senior equity claim with limited rights & Assuming bond-equivalent payment certainty \\
Liquidating value & Estimated net proceeds distributable in winding down & Ignoring costs, liens and enforcement delays \\
Book value & Accounting net asset position & Assuming resale proceeds equal carrying value \\
Capitalization structure & Distribution of claims on business cash and assets & Focusing only on nominal share price \\
Price--value discrepancy & Difference between quote and underwriting value & Assuming a low quote reveals mispricing \\
\bottomrule
\end{longtable}
\normalsize

\section{Engineering acceptance criteria}
A credible investment analysis must (i) identify the security and actual obligor, (ii) tie every financial quantity to a date and source, (iii) separately show bear/base/bull cases, (iv) state how a permanent loss occurs, (v) document valuation sensitivity and liquidity constraints, (vi) receive independent challenge, and (vii) retain a versioned decision and monitoring plan.

\subsection{Selected book-navigation map}
\begin{center}\small\begin{tabularx}{\textwidth}{p{0.12\textwidth}X X}\toprule Part & Historical area & Engineering application \\
\midrule
I & Survey, classification, sources & Evidence controls and investment charter \\
II & Fixed-value investments & Credit underwriting and covenant risk \\
III & Speculative senior securities & Embedded optionality and distress \\
IV & Common stocks and dividends & Earning-power and payout evaluation \\
V & Income account & Normalization and dilution \\
VI & Balance-sheet analysis & Net asset, liquidation, hidden liabilities \\
VII & Additional price--value aspects & Comparables, governance, mispricing \\
VIII & Contemporary viewpoints & International, private and endowment contexts \\
\bottomrule\end{tabularx}\end{center}

\end{document}
